\documentclass[11pt,a4paper]{article}

\usepackage{amsmath,amssymb,amsthm}
\usepackage{mathtools}
\usepackage{hyperref}
\usepackage{cleveref}
\usepackage[margin=1in]{geometry}
\usepackage{enumitem}
\usepackage{booktabs}

\newtheorem{theorem}{Theorem}[section]
\newtheorem{lemma}[theorem]{Lemma}
\newtheorem{proposition}[theorem]{Proposition}
\newtheorem{corollary}[theorem]{Corollary}
\newtheorem{definition}[theorem]{Definition}
\newtheorem{axiom}{Axiom}
\newtheorem{remark}[theorem]{Remark}
\newtheorem{example}[theorem]{Example}

\newcommand{\Prob}{\mathsf{Prob}}
\newcommand{\ProbProp}{\mathsf{ProbProp}}
\newcommand{\ProbNat}{\mathsf{ProbNat}}
\newcommand{\E}{\mathbb{E}}
\newcommand{\pand}{\wedge_p}
\newcommand{\por}{\vee_p}
\newcommand{\pnot}{\neg_p}
\newcommand{\pzero}{\mathbb{0}}
\newcommand{\pone}{\mathbb{1}}
\newcommand{\ple}{\leq_p}
\newcommand{\padd}{+_p}
\newcommand{\pmul}{\cdot_p}
\newcommand{\psub}{-_p}
\newcommand{\psum}{\Sigma_p}
\newcommand{\Nat}{\mathbb{N}}

\title{Probabilistic Natural Numbers\\[0.5em]
\large Distributions over $\Nat$ as Mathematical Objects}

\author{Probabilistic Foundations Project}

\date{\today}

\begin{document}

\maketitle

\begin{abstract}
We develop natural numbers within the probabilistic foundations framework. A probabilistic natural number is not a single value but a distribution over $\Nat$. The Peano axioms hold in this setting: zero is not a successor, and the successor operation shifts distributions. Deterministic naturals (point masses) embed as a special case, but the full framework supports reasoning about uncertain quantities, random variables, and probabilistic bounds. This tests the expressiveness of probabilistic foundations: can we do number theory without assuming numbers are determinate?
\end{abstract}

\section{Introduction}

\subsection{Why Probabilistic Natural Numbers?}

In classical mathematics, natural numbers are determinate: the number 7 is exactly 7. But many mathematical situations involve uncertain quantities:

\begin{itemize}
    \item The number of prime factors of a random integer
    \item The runtime of a randomized algorithm
    \item The population size after stochastic dynamics
    \item Counting under measurement uncertainty
\end{itemize}

Rather than building classical naturals and then adding probability, we build \emph{probabilistic naturals directly}. A $\ProbNat$ is a distribution over possible values.

\subsection{Deterministic Naturals as Special Case}

The deterministic natural number $n$ is the point mass at $n$: probability 1 of being $n$, probability 0 of being anything else. Classical number theory embeds in this framework.

But the interesting work is with genuine distributions: geometric distributions, Poisson distributions, or arbitrary discrete distributions over $\Nat$.

\subsection{Prerequisites}

This paper assumes Papers 1--2: the $\Prob$ type, $\ProbProp$, and conditional expectation.

\section{Countable Summation}

\subsection{A New Primitive}

To define distributions over $\Nat$, we need to sum countably many probabilities.

\begin{axiom}[Countable Summation]
There exists an operation:
\[
\psum : (\Nat \to \Prob) \to \Prob
\]
\end{axiom}

\begin{axiom}[Singleton Sum]
$\psum\left(\lambda m.\, \begin{cases} p & m = n \\ \pzero & m \neq n \end{cases}\right) = p$
\end{axiom}

\begin{axiom}[Zero Sum]
$\psum(\lambda n.\, \pzero) = \pzero$
\end{axiom}

\begin{axiom}[Bounded Sum]
If $f(n) \ple \pone$ for all $n$, then $\psum(f) \ple \pone$.
\end{axiom}

\begin{remark}
We do not require countable additivity as an axiom; it would follow from a construction of $\Prob$ in Paper 5.
\end{remark}

\section{Probabilistic Natural Numbers}

\subsection{Definition}

\begin{definition}[ProbNat]
A \emph{probabilistic natural number} over type $\alpha$ is a structure:
\[
N : \ProbNat(\alpha)
\]
consisting of:
\begin{enumerate}
    \item A family $\mathsf{is\_n} : \Nat \to \ProbProp(\alpha)$ where $N.\mathsf{is\_n}(k)(x)$ is the probability that $N$ equals $k$ at state $x$
    \item Exhaustiveness: $\psum(\lambda k.\, N.\mathsf{is\_n}(k)(x)) = \pone$ for all $x$
    \item Disjointness: $N.\mathsf{is\_n}(k)(x) \pmul N.\mathsf{is\_n}(m)(x) = \pzero$ when $k \neq m$
\end{enumerate}
\end{definition}

\begin{remark}
Exhaustiveness says: the probabilities sum to 1 (it's a distribution).
Disjointness says: you can't be two different values simultaneously.
\end{remark}

\subsection{Interpretation}

A $\ProbNat$ is a random variable taking values in $\Nat$. At each state $x : \alpha$, it specifies a probability distribution over possible values.

\begin{example}[Fair Die]
A fair six-sided die as $\ProbNat$:
\[
D.\mathsf{is\_n}(k)(x) = \begin{cases} 1/6 & 1 \leq k \leq 6 \\ 0 & \text{otherwise} \end{cases}
\]
\end{example}

\begin{example}[Geometric Distribution]
Number of coin flips until first heads (geometric with $p = 0.5$):
\[
G.\mathsf{is\_n}(k)(x) = \begin{cases} 0.5^k & k \geq 1 \\ 0 & k = 0 \end{cases}
\]
\end{example}

\section{Zero and Successor}

\subsection{Zero}

\begin{definition}[Zero Proposition]
For $N : \ProbNat(\alpha)$:
\[
\mathsf{zero}(N) := N.\mathsf{is\_n}(0)
\]
The probabilistic proposition ``$N$ equals zero.''
\end{definition}

\subsection{Successor}

The successor operation shifts a distribution: $\mathsf{succ}(N)$ has the same shape as $N$ but shifted by one.

\begin{definition}[Successor Is-N]
For $N : \ProbNat(\alpha)$:
\[
\mathsf{succ\_is\_n}(N, k) := \begin{cases}
    \mathbf{0} & k = 0 \\
    N.\mathsf{is\_n}(k-1) & k > 0
\end{cases}
\]
where $\mathbf{0}(x) = \pzero$ for all $x$.
\end{definition}

\begin{theorem}[Successor Preserves Exhaustiveness]
If $N$ is exhaustive, so is $\mathsf{succ\_is\_n}(N, \cdot)$.
\end{theorem}

\begin{proof}
\begin{align*}
    \psum(\lambda k.\, \mathsf{succ\_is\_n}(N, k)(x))
    &= \psum\left(\lambda k.\, \begin{cases} \pzero & k = 0 \\ N.\mathsf{is\_n}(k-1)(x) & k > 0 \end{cases}\right) \\
    &= \psum(\lambda m.\, N.\mathsf{is\_n}(m)(x)) \quad \text{(reindexing)} \\
    &= \pone
\end{align*}
\end{proof}

\begin{theorem}[Successor Preserves Disjointness]
If $N$ is disjoint, so is $\mathsf{succ\_is\_n}(N, \cdot)$.
\end{theorem}

\begin{proof}
For $k \neq m$:
\begin{itemize}
    \item If $k = 0$: $\mathsf{succ\_is\_n}(N, 0)(x) = \pzero$, so the product is $\pzero$.
    \item If $m = 0$: Similarly.
    \item If $k, m > 0$: $\mathsf{succ\_is\_n}(N, k) \pmul \mathsf{succ\_is\_n}(N, m) = N.\mathsf{is\_n}(k-1) \pmul N.\mathsf{is\_n}(m-1) = \pzero$ since $k-1 \neq m-1$.
\end{itemize}
\end{proof}

\section{The Peano Axioms}

\subsection{Peano 1: Zero is Not a Successor}

\begin{theorem}[Peano 1]
For any $N : \ProbNat(\alpha)$ and $x : \alpha$:
\[
\mathsf{succ\_is\_n}(N, 0)(x) = \pzero
\]
\end{theorem}

\begin{proof}
By definition: $\mathsf{succ\_is\_n}(N, 0) = \mathbf{0}$.
\end{proof}

\begin{remark}
This says: the successor of any distribution has zero probability of being zero. You can't ``shift down'' past zero.
\end{remark}

\subsection{Peano 2: Successor is Injective}

Classical Peano 2 says: if $\mathsf{succ}(n) = \mathsf{succ}(m)$ then $n = m$.

In our setting, ``equality'' of distributions means: same probability of each value.

\begin{theorem}[Peano 2 / Shift Property]
$\mathsf{succ\_is\_n}(N, k+1)(x) = N.\mathsf{is\_n}(k)(x)$
\end{theorem}

\begin{proof}
By definition: $\mathsf{succ\_is\_n}(N, k+1) = N.\mathsf{is\_n}(k)$.
\end{proof}

\begin{remark}
This captures injectivity: the successor operation is a bijective shift on distributions. The probability that $\mathsf{succ}(N) = k+1$ equals the probability that $N = k$.
\end{remark}

\subsection{Induction (Sketch)}

Probabilistic induction says: if $P$ holds for the deterministic zero with probability 1, and $P(\mathsf{succ}(N))$ follows from $P(N)$ for all $N$, then $P$ holds for all $\ProbNat$.

The full formulation requires quantification over distributions, which we defer to Paper 5.

\section{Deterministic Natural Numbers}

\subsection{Point Masses}

\begin{definition}[Deterministic Natural]
The deterministic natural number $n$ is:
\[
\mathsf{det}(n).\mathsf{is\_n}(k)(x) := \begin{cases} \pone & k = n \\ \pzero & k \neq n \end{cases}
\]
\end{definition}

\begin{theorem}[Det is ProbNat]
$\mathsf{det}(n)$ satisfies exhaustiveness and disjointness.
\end{theorem}

\begin{proof}
Exhaustiveness: $\psum$ of a singleton is $\pone$ (Singleton Sum axiom).
Disjointness: at most one value is non-zero.
\end{proof}

\subsection{Classical Naturals Embed}

\begin{theorem}
$\mathsf{det}(0)$ is the unique distribution with $\mathsf{is\_n}(0) = \mathbf{1}$.
\end{theorem}

\begin{theorem}[Successor of Deterministic]
$\mathsf{succ\_is\_n}(\mathsf{det}(n), n+1)(x) = \pone$
\end{theorem}

\begin{proof}
$\mathsf{succ\_is\_n}(\mathsf{det}(n), n+1) = \mathsf{det}(n).\mathsf{is\_n}(n) = \pone$.
\end{proof}

Thus the successor of the deterministic $n$ is (with probability 1) the value $n+1$.

\section{Arithmetic on Distributions}

\subsection{Expected Value}

For a $\ProbNat$, the expected value is:
\[
\E[N](x) := \psum(\lambda k.\, k \cdot N.\mathsf{is\_n}(k)(x))
\]

(This requires extending $\Prob$ operations to handle naturals, or treating expectation as a derived quantity.)

\subsection{Addition of Independent ProbNats}

If $M, N$ are independent $\ProbNat$s, their sum has:
\[
(M + N).\mathsf{is\_n}(k)(x) = \psum\left(\lambda j.\, M.\mathsf{is\_n}(j)(x) \pmul N.\mathsf{is\_n}(k-j)(x)\right)
\]

This is convolution of distributions.

\subsection{Comparison}

\begin{definition}[Probabilistic Less-Than]
$P(M < N)(x) := \psum\left(\lambda k.\, M.\mathsf{is\_n}(k)(x) \pmul \psum\left(\lambda m.\, N.\mathsf{is\_n}(m)(x) \cdot [m > k]\right)\right)$
\end{definition}

where $[m > k] = \pone$ if $m > k$, else $\pzero$.

\section{Examples}

\subsection{Geometric Distribution}

Let $G$ be geometric with parameter $p$:
\[
G.\mathsf{is\_n}(k)(x) = (1-p)^k \cdot p \quad (k \geq 0)
\]

This represents the number of failures before first success.

\begin{itemize}
    \item $\E[G] = (1-p)/p$
    \item $\mathsf{succ}(G)$ shifts: now counts from 1, not 0
    \item $G + G'$ (independent) is negative binomial
\end{itemize}

\subsection{Poisson Distribution}

Let $P_\lambda$ be Poisson with rate $\lambda$:
\[
P_\lambda.\mathsf{is\_n}(k)(x) = e^{-\lambda} \cdot \lambda^k / k!
\]

This models event counts in fixed time intervals.

\subsection{Uniform on Finite Set}

Uniform distribution on $\{0, 1, \ldots, n-1\}$:
\[
U_n.\mathsf{is\_n}(k)(x) = \begin{cases} 1/n & k < n \\ 0 & k \geq n \end{cases}
\]

\section{Why Not Just Use Random Variables?}

\subsection{The Classical Approach}

In classical mathematics:
\begin{enumerate}
    \item Define $\Nat$ inductively
    \item Build measure theory on $\mathbb{R}$ (or general spaces)
    \item Define random variables as measurable functions
    \item A ``random natural'' is a random variable $X : \Omega \to \Nat$
\end{enumerate}

\subsection{The Probabilistic Foundations Approach}

We invert this:
\begin{enumerate}
    \item Define $\ProbNat$ directly as distributions
    \item Deterministic naturals are the $\{0,1\}$-valued special case
    \item No measure theory prerequisites
\end{enumerate}

\subsection{Advantages}

\begin{itemize}
    \item Uncertainty is primitive, not derived
    \item Reasoning about distributions is native
    \item No probability space $\Omega$ is needed (distributions are first-class)
    \item Foundation is more direct for probabilistic applications
\end{itemize}

\section{Implementation}

\subsection{Lean 4 Code}

\begin{verbatim}
axiom prob_sum : (Nat -> Prob) -> Prob
axiom prob_sum_singleton : forall n p,
  prob_sum (fun m => if m = n then p else 0) = p

structure ProbNat (a : Type) where
  is_n : Nat -> ProbProp a
  exhaustive : forall x, prob_sum (fun n => is_n n x) = 1
  disjoint : forall n m, n <> m -> forall x, is_n n x * is_n m x = 0

noncomputable def prob_succ_is_n (N : ProbNat a) (n : Nat) : ProbProp a :=
  match n with
  | 0 => fun _ => 0
  | Nat.succ m => N.is_n m

theorem peano1 (N : ProbNat a) (x : a) :
    prob_succ_is_n N 0 x = 0 := rfl

theorem peano2_shift (N : ProbNat a) (k : Nat) (x : a) :
    prob_succ_is_n N (Nat.succ k) x = N.is_n k x := rfl

noncomputable def det_nat (n : Nat) : ProbNat a where
  is_n := fun m => fun _ => if m = n then 1 else 0
  exhaustive := by intro x; exact prob_sum_singleton n 1
  disjoint := by
    intro i j hij x
    by_cases hi : i = n <;> by_cases hj : j = n
    . exact absurd (hi.trans hj.symm) hij
    . simp [hi, hj]
    . simp [hi, hj]
    . simp [hi, hj]
\end{verbatim}

\subsection{Verification Status}

All Peano axioms and closure properties are machine-verified.

\section{Discussion}

\subsection{What This Demonstrates}

We have shown that natural numbers can be developed within probabilistic foundations:
\begin{itemize}
    \item The Peano axioms hold
    \item Deterministic naturals embed faithfully
    \item Genuine distributions are first-class objects
\end{itemize}

\subsection{Limitations}

The current development is axiomatic. A synthetic construction of $\psum$ from more primitive operations would strengthen the foundations (Paper 5).

\subsection{What Comes Next}

\begin{itemize}
    \item \textbf{Paper 4}: Using probabilistic methods to prove existence (Zorn without AC)
    \item \textbf{Paper 5}: Synthetic construction of $\Prob$
\end{itemize}

\section{Conclusion}

Probabilistic natural numbers are distributions over $\Nat$. The Peano axioms hold: zero is not a successor, successor is injective (shifts distributions). Deterministic naturals are the point-mass special case.

This demonstrates that number theory can be developed within probabilistic foundations---not by first assuming classical naturals and then adding probability, but by building probabilistic naturals directly.

\bibliographystyle{plain}

\end{document}
